% Generated by scripts/assemble.py; corrections belong in source/corrections.json.
\examhead{102}{浙江大学 2025 年数学分析试卷}{101}
% p102-b20; source PDF page 102
\par\Needspace{5\baselineskip}
\blockanchor{p102-b20}
\majorquestion{一}计算题.\par
% p102-b21; source PDF page 102
\blockanchor{p102-b21}
\question{1}$\lim_{n\to \infty}n^2\left(\arctan{\frac{2024}{n}} -\arctan{\frac{2024}{n + 1}}\right).$\par
% p102-b22; source PDF page 102
\blockanchor{p102-b22}
\question{2}$z = z(x,y)$ 由 $f(z - x,z - y)$ 确定，$f(u,v)$ 有一阶偏导数，且 $\frac{\partial f}{\partial u} + \frac{\partial f}{\partial v} \neq 0$，计算 $\frac{\partial z}{\partial x} + \frac{\partial z}{\partial y}$.\par
% p102-b23; source PDF page 102
\blockanchor{p102-b23}
\question{3}计算 $\int_{-8}^{8} \frac{1 + \sin x}{1 + \sqrt[3]{x^2}} \, \mathrm{d}x$.\par
% p102-b24; source PDF page 102
\par\Needspace{6\baselineskip}
\blockanchor{p102-b24}
\question{4}$\Sigma = \{(x, y, z) \mid x^2 + y^2 + z^2 = 1, z \geq 0\}$, 取上侧, 计算\par
\nopagebreak[4]
% p102-b25; source PDF page 102
\blockanchor{p102-b25}
\displayformula{\iint_ {\Sigma} x ^ {2} \mathrm{d} y \mathrm{d} z + y ^ {2} \mathrm{d} z \mathrm{d} x + [ z ^ {2} + \sin (x y) ] \mathrm{d} x \mathrm{d} y.}
% p103-b00; source PDF page 103
\par\Needspace{5\baselineskip}
\blockanchor{p103-b00}
\majorquestion{二}设 $\{a_{n}\}$ 为非负递减数列, 且级数 $\sum_{n=1}^{\infty} a_{n}$ 收敛, 证明: $\lim_{n \to \infty} na_{n} = 0$.\par
% p103-b01; source PDF page 103
\par\Needspace{5\baselineskip}
\blockanchor{p103-b01}
\majorquestion{三}设 $f(x)$ 在闭区间 $[0,1]$ 上连续, 且对任意的 $(0,1)$, $f_{+}^{\prime}(x)$ 存在. 若 $f(0) = f(1)$, 证明: 存在 $\xi \in (0,1)$, 使得 $f_{+}^{\prime}(\xi) \leq 0$.\par
% p103-b02; source PDF page 103
\par\Needspace{5\baselineskip}
\blockanchor{p103-b02}
\majorquestion{四}叙述闭区间套定理, 用该定理证明闭区间上连续函数可达最大值.\par
% p103-b03; source PDF page 103
\par\Needspace{5\baselineskip}
\par\Needspace{6\baselineskip}
\blockanchor{p103-b03}
\majorquestion{五}设\par
\nopagebreak[4]
% p103-b04; source PDF page 103
\blockanchor{p103-b04}
\displayformula{f (x, y) = \left\{ \begin{array}{l l} (x + y) ^ {2} \sin \frac {1}{x ^ {2} + y ^ {2}}, & (x, y) \neq (0, 0); \\ 0, & (x, y) = (0, 0). \end{array} \right.}
% p103-b05; source PDF page 103
\blockanchor{p103-b05}
讨论 f 和 f 的偏导数在  $(0,0)$  处的连续性及 f 在  $(0,0)$  的可微性.\par
% p103-b06; source PDF page 103
\par\Needspace{5\baselineskip}
\blockanchor{p103-b06}
\majorquestion{六}证明: $\sum_{n=1}^{\infty} n^{2} e^{-nx}$ 在 $(0, +\infty)$ 上非一致收敛, 在 $[1, +\infty)$ 上一致收敛.\par
% p103-b07; source PDF page 103
\par\Needspace{5\baselineskip}
\blockanchor{p103-b07}
\majorquestion{七}证明: $f(x)$ 在 $(0,1)$ 上一致连续的充要条件是对任意 $(0,1)$ 上的 Cauchy 列 $\{x_{n}\}, \{f(x_{n})\}$ 也为 Cauchy 列.\par
% p103-b08; source PDF page 103
\par\Needspace{5\baselineskip}
\par\Needspace{6\baselineskip}
\blockanchor{p103-b08}
\majorquestion{八}设 $\{r_n\}_{n=1}^{\infty}$ 为全体有理数, 令\par
\nopagebreak[4]
% p103-b09; source PDF page 103
\blockanchor{p103-b09}
\displayformula{f (x) = \sum_ {n: r _ {n} <   x} \frac {1}{2 ^ {n}}.}
% p103-b10; source PDF page 103
\blockanchor{p103-b10}
其中上述和式是对所有满足 $r_n < x$ 的正整数 $n$ 求和. 证明: $f(x)$ 在无理点连续, 在有理点处间断, 且单调递增.\par
% p103-b11; source PDF page 103
\par\Needspace{5\baselineskip}
\par\Needspace{6\baselineskip}
\blockanchor{p103-b11}
\majorquestion{九}对任意的实数 $x$ 与正整数 $n$, 记 $\varphi_n(x) = \frac{n}{\sqrt{\pi}} e^{-n^2 x^2}$. 证明: 对任意在实轴上连续且有界的函数 $f$, 对任意的实数 $x$, 有\par
\nopagebreak[4]
% p103-b12; source PDF page 103
\blockanchor{p103-b12}
\displayformula{\lim _ {n \rightarrow \infty} \int_ {- \infty} ^ {+ \infty} f (y) \varphi_ {n} (x - y) d y = f (x).}
